% =============================================================
%  Preamble — BUBT CSE Capstone formatting
% =============================================================

% Page geometry: a strict 1in margin from page edge to body text on ALL four sides.
% The header is empty (\fancyhf{} below) so we set headheight and headsep to 0,
% meaning the body literally starts 1in from the top edge — no invisible header
% reservation pushing it down. footskip places the page-number footer inside the
% 1in bottom margin.
\usepackage[a4paper,
            top=1in,bottom=1in,left=1in,right=1in,
            headheight=0pt,headsep=0pt,footskip=20pt]{geometry}

% Times-compatible font stack.
% newtxtext provides true bold weight (700) for headings; newtxmath provides
% Times-matched math. Loading both keeps text and math in the same family,
% eliminating mismatched-font warnings from mathptmx (zptmcm*, ptmb7t, ptmri7t,
% msam10, msbm10) and the lighter-bold fallback the old setup had.
\usepackage[T1]{fontenc}
\usepackage{amsmath}        % must load BEFORE newtxmath
\usepackage{newtxtext}
\usepackage{newtxmath}      % Times-compatible math, also provides \mathbb, \Bbbk, etc.

% 1.5 line spacing
\usepackage{setspace}
\setstretch{1.15}

% Justified text
\usepackage{ragged2e}

% Headings
\usepackage[explicit]{titlesec}
\usepackage{titletoc}

% Allow up to 4 levels of section numbering (1.1.1.1)
\setcounter{secnumdepth}{4}
\setcounter{tocdepth}{3}

% Numbered chapter heading (Ch 1–6, appendix): "Chapter N" centered huge bold,
% then "N. Title" left-aligned bold (matches BUBT example reports).
% The \vspace*{-0.5in} pulls the heading up so its top sits closer to the page
% edge — visually balancing the top space with the side margins.
\titleformat{\chapter}[display]
  {\vspace*{-0.5in}\normalfont\bfseries}
  {\filcenter\huge\bfseries Chapter \thechapter}
  {24pt}
  {\LARGE\bfseries \thechapter.\enspace #1}
\titlespacing*{\chapter}{0pt}{0pt}{18pt}

% Unnumbered (starred) chapter heading — used by frontmatter (abstract,
% dedication, etc.) and references. Just the title, centered, no "Chapter N"
% line and no leading number. Same negative top-shift as the numbered version.
\titleformat{name=\chapter,numberless}[display]
  {\vspace*{-0.5in}\normalfont\bfseries}
  {}
  {0pt}
  {\filcenter\huge\bfseries #1}
\titlespacing*{name=\chapter,numberless}{0pt}{0pt}{18pt}

% Hierarchy: section > subsection > subsubsection in size; all bold.
\titleformat{\section}{\normalfont\Large\bfseries}{\thesection.}{0.6em}{#1}
\titleformat{\subsection}{\normalfont\large\bfseries}{\thesubsection.}{0.6em}{#1}
\titleformat{\subsubsection}{\normalfont\normalsize\bfseries}{\thesubsubsection.}{0.6em}{#1}
\titleformat{\paragraph}[runin]{\normalfont\normalsize\bfseries}{}{0pt}{#1}

% Tightened section/subsection/subsubsection spacing (S4 of tightening pass)
\titlespacing*{\section}{0pt}{8pt plus 2pt minus 1pt}{4pt}
\titlespacing*{\subsection}{0pt}{6pt plus 2pt minus 1pt}{3pt}
\titlespacing*{\subsubsection}{0pt}{5pt plus 1pt minus 1pt}{2pt}

% Block-style paragraphs (no first-line indent + vertical gap between),
% matching the BUBT example reports.
\setlength{\parindent}{0pt}
\setlength{\parskip}{0.35\baselineskip plus 1pt minus 1pt}

% Headers/footers — page number bottom-right only
\usepackage{fancyhdr}
\pagestyle{fancy}
\fancyhf{}
\fancyfoot[R]{\thepage}
\renewcommand{\headrulewidth}{0pt}
\renewcommand{\footrulewidth}{0pt}
\fancypagestyle{plain}{
  \fancyhf{}
  \fancyfoot[R]{\thepage}
  \renewcommand{\headrulewidth}{0pt}
}

% Graphics
\usepackage{graphicx}
\graphicspath{{figures/}}

% Tables
\usepackage{booktabs}
\usepackage{array}
\usepackage{longtable}
\usepackage{multirow}
\usepackage{makecell}
\usepackage{tabularx}
\usepackage{float}   % provides [H] to force float placement

% placeins: provides \FloatBarrier to prevent floats from drifting past a
% specified boundary. With [section] option, the barrier is applied
% automatically at every \section so a figure declared in §4.2 cannot
% physically appear inside §4.3's text, which would otherwise split a
% sentence across the figure.
\usepackage[section]{placeins}
% Extend the barrier to subsection and subsubsection too: a figure declared
% inside a subsection can never physically appear inside the next subsection
% or subsubsection. This is what prevents sentence-splitting by drifting floats.
\usepackage{etoolbox}
\preto\subsection{\FloatBarrier}
\preto\subsubsection{\FloatBarrier}

% Captions: bold "Figure N:" / "Table N:" then descriptive text
\usepackage[labelfont=bf,labelsep=colon,format=plain,
            justification=centering,font=small]{caption}
\usepackage{subcaption}

% Tighten whitespace around floats + captions (page-reduction + tightening passes)
\setlength{\textfloatsep}{5pt plus 1pt minus 1pt}
\setlength{\floatsep}{5pt plus 1pt minus 1pt}
\setlength{\intextsep}{5pt plus 1pt minus 1pt}
\setlength{\abovecaptionskip}{10pt}
\setlength{\belowcaptionskip}{0pt}

% Anchor floats to the TOP of float pages (default \@fptop has 1fil
% vertical glue, which centers the float vertically). With these settings,
% any figure that LaTeX puts on its own page lands flush with the top margin,
% subsequent floats are stacked with a tight 10pt gap (no stretch), and any
% leftover space accumulates at the bottom.
\makeatletter
\setlength{\@fptop}{0pt}
\setlength{\@fpsep}{10pt}
\setlength{\@fpbot}{0pt plus 1fil}
\makeatother

% Tighter display-equation spacing (S5 of tightening pass)
\setlength{\abovedisplayskip}{4pt plus 1pt minus 1pt}
\setlength{\belowdisplayskip}{4pt plus 1pt minus 1pt}
\setlength{\abovedisplayshortskip}{2pt plus 1pt minus 1pt}
\setlength{\belowdisplayshortskip}{2pt plus 1pt minus 1pt}

% Default array stretch (S1 of tightening pass): override per-table 1.05 settings.
\renewcommand{\arraystretch}{1.0}

% Math
% amsmath already loaded at top (before newtxmath). amssymb dropped:
% newtxmath provides \mathbb, \Bbbk, and the rest. Re-adding amssymb would
% conflict with newtxmath's \Bbbk definition.

% Lists
\usepackage{enumitem}
% Keep lists compact so global \parskip doesn't inflate item spacing.
\setlist{topsep=0.15\baselineskip, partopsep=0pt, itemsep=0pt, parsep=0pt}

% IEEE-style numeric citations: compresses [1,2,3] into [1--3]
% Loaded BEFORE hyperref per the cite package documentation.
\usepackage{cite}

% Hyperlinks
\usepackage[colorlinks=false,hidelinks]{hyperref}

% References — using bibtex (built into Tectonic, no extra install needed)
% IEEEtran.bst gives IEEE numeric style
% Rename the auto-generated bibliography heading from "Bibliography" to "References"
\renewcommand{\bibname}{References}

% URL formatting (xurl breaks long URLs at any character to prevent overfull)
\usepackage{xurl}
\Urlmuskip=0mu plus 1mu

% Microtype: protrusion + expansion → fewer overfull/underfull warnings
\usepackage{microtype}

% Allow more flexible word/line breaks for borderline-fitting paragraphs
\sloppy
\emergencystretch=5em
\hbadness=10000   % suppress underfull hbox warnings (cosmetic only)
\hfuzz=2pt        % tolerate minor (<2pt) overfull boxes silently

% Landscape pages for wide tables
\usepackage{pdflscape}
\usepackage{lscape}

% Avoid orphaned headings
\usepackage[bottom]{footmisc}
\widowpenalty=10000
\clubpenalty=10000
% Forbid a hyphenated word from being split across a page break
% (the second half of "de- scribe" must not land on the next page).
\brokenpenalty=10000

% Custom title-page rule (renamed to avoid clash with titlesec's \titlerule)
\newcommand{\titlepagerule}{\noindent\rule{\textwidth}{0.5pt}}

% Chapter-prefixed equation numbering (3.1, 3.2, ...). Numbering is
% fully automatic: each \begin{equation} block gets the next number,
% and equations renumber automatically if any are added, removed, or
% reordered. Reference any equation with \eqnref{eq:label} → "Eq. 3.4".
\numberwithin{equation}{chapter}
\newcommand{\eqnref}[1]{Eq.~\ref{#1}}
\newcommand{\Eqnref}[1]{Equation~\ref{#1}}
